\documentclass[10pt,a4paper]{amsart}
\usepackage[margin=19mm]{geometry}
\usepackage{amsmath,amssymb,mathtools}
\usepackage[scr=rsfs,cal=euler]{mathalfa}
\usepackage{xfrac}
\usepackage[unicode,hidelinks,bookmarks=false]{hyperref}
\newcommand{\ie}{i.e.\ }
\newcommand{\cf}{cf.\ }
\newcommand{\resp}{respectively\ }
\newcommand{\Subsection}{\subsection}
\providecommand{\nobreakdash}{\nobreak\hbox{-}}
\setlength{\emergencystretch}{2em}
\multlinegap0pt
\usepackage{bm}
\usepackage{bbm}
\usepackage{dsfont}
\usepackage{xspace}
\usepackage{cancel}
\usepackage{mathtools}
\usepackage{amsthm}
\makeatletter
\@ifundefined{theorem}{\newtheorem{theorem}{Theorem}}{}
\makeatother
\newcommand{\B}{\mathcal{B}}
\newcommand{\K}{\mathcal K}
\renewcommand{\H}{\mathcal{H}}
\title[Representations of noncommutative cubes and prisms]{Representations of noncommutative cubes and prisms}
\author{Douglas Farenick, Roghayeh Maleki, Sofia Medina Varela,, Sushil Singla}
\date{Source: arXiv:2601.16902v2 (CC BY 4.0)}
\begin{document}
\maketitle

\noindent\emph{Source and licence.} Representations of noncommutative cubes and prisms. Version arXiv:2601.16902v2 (CC BY 4.0), distributed under the Creative Commons Attribution 4.0 International licence, as recorded in the arXiv licence metadata for this exact version and shown on the abstract page https://arxiv.org/abs/2601.16902v2. Full rights notice: licenses/arxiv-2601.16902-CC-BY-4.0.txt.\par\smallskip

\noindent\textit{Editorial scope.} Counted unit: the Theorem whose source label is \texttt{thm d=2} in the article (source0.tex lines 643--646, source label \texttt{thm d=2}), delivered together with the article's own complete original proof of it (source0.tex lines 648--695, one \texttt{proof} environment of 48 lines). No mathematics is written, repaired or re-derived: statement and proof are the article's own text line for line. Required context retained uncounted: none. Every definition, display and label of those lines is retained unchanged. Omitted from this package and not redistributed: 1--642, 647--647, 696--2066. Nothing inside a retained excerpt is omitted or altered: every retained body is delivered exactly as the article's lines are written, so no omission receipt of any kind accompanies this package. Citations: the retained excerpts carry no citation command at all, so no bibliography entry of the article is transcribed and no reference is redistributed. Reference adaptation: none was needed and none was made; every reference inside the delivered unit resolves inside it, so no unresolved reference or citation is produced. Printed slips kept literal: no printed word, symbol, hypothesis or number of the counted statement or of its proof was altered. Layout and class: the article's own class is \texttt{amsart} with packages including amscd, latexsym, amssymb, palatino, color, graphicx, amsmath, amsthm; the wrapper is the lane's pinned \texttt{amsart} editorial wrapper at 10pt on A4, which restates the theorem-like environments the retained text needs and adds the article's own macro definitions that the retained text needs (\texttt{\textbackslash B}, \texttt{\textbackslash H}, \texttt{\textbackslash K}), with extra wrapper packages (bm, bbm, dsfont, xspace, cancel, mathtools). The delivered numbering is the unit's own. Assets: no retained span contains a figure environment, a graphics-inclusion command, a PDF-page inclusion, a table or a TikZ picture; no image file is copied into this package, the package declares no asset dependency and no image file is redistributed with it. Licence: version arXiv:2601.16902v2 of this article is distributed under the Creative Commons Attribution 4.0 International licence, as recorded in the arXiv licence metadata for this exact version and shown on the abstract page https://arxiv.org/abs/2601.16902v2. A full rights notice is delivered at licenses/arxiv-2601.16902-CC-BY-4.0.txt. Characters: every retained excerpt is pure ASCII, so the delivered engine, XeLaTeX with its default Latin Modern text font, reports no missing character.
\begin{theorem}\label{thm d=2}
Assume $v_1,v_2\in\B(\H)$ are linearly independent symmetries.  
	If $\{v_1,v_2\}$ is irreducible, then $\mbox{\textrm dim}\H=2$.
\end{theorem}

\begin{proof} Decompose $\H$ as $\K_{-1}\oplus\K_{1}$, 
	where $\K_j$ is the eigenspace of $v_j$ corresponding to the eigenvalue $j\in\{-1,1\}$. Thus, $v_1$ and
	$v_2$ have the following representations as $2\times 2$ operator matrices acting on $\H=\K_{-1}\oplus\K_{1}$:
	\[
	v_1=\left[ \begin{array}{cc} -1_{\K_{-1}} &0 \\ 0 & 1_{\K_1} \end{array}\right]
	\,\mbox{ and }\,
	v_2=\left[ \begin{array}{cc} a_{-1}&y \\ y^* & a_{1} \end{array}\right],
	\]
	for some selfadjoint $a_j\in\B(\K_j)$ and bounded linear operator $y:\K_{1}\rightarrow\K_{-1}$.
	We now determine the elements $a_j$ and $y$, using the hypothesis that $\{v_1,v_2\}$ is an irreducible set.
	
	Every matrix that commutes with $v_1$ has the form $x_{-1}\oplus x_{1}$, for $x_j\in\B(\K_j)$.
	Because $v_2$ is a symmetry, $ya_{1}=-a_{-1}y$. Hence, the diagonal matrix 
	$-a_{-1}\oplus a_{1}$
	commutes with both $v_1$ and $v_2$, and so there exists
	$\lambda\in\mathbb R$ such that $a_{-1}=\lambda1_{\K_{-1}}$ and
	$a_1=-\lambda1_{\K_1}$. The diagonal entries of the matrix 
	$v_2^2=1_\H$ lead to the equations $\lambda^2 1_{\K_{-1}}+yy^*=1_{\K_{-1}}$ and
	$\lambda^2 1_{\K_{1}}+y^*y=1_{\K_{1}}$. Thus, $-1\leq\lambda\leq1$. However, $|\lambda|\not=1$, 
	for if it were true that $|\lambda|=1$, then we would necessarily have $y=0$, implying $v_2=\lambda v_1$,
	which is in contradiction to the linear independence of $v_1$ and $v_2$. Hence, $\lambda\in(-1,1)$.
	
	Since $yy^*=(1-\lambda^2)\,1_{\K_{-1}}$ and $y^*y=(1-\lambda^2)\,1_{\K_{1}}$, the element $z=\frac{1}{\sqrt{1-\lambda^2}}y$ 
	is a linear operator $z:\K_{1}\rightarrow\K_{{-1}}$ such that $zz^*=1_{\K_{-1}}$ and $z^*z=1_{\K_{1}}$; that is, $z$ is both an isometry
	and a co-isometry. Hence, $z:\K_{1}\rightarrow\K_{-1}$ is a surjective isometry, and we can assume, henceforth
	and without loss of generality, that $\K_{-1}=\K_1$, which we denote by $\K$. Thus, $v_2$ is given by
	\[
	v_2=\left[\begin{array}{cc} \lambda 1_\K & y \\ y^* & -\lambda 1_\K\end{array}\right],
	\]
	where $yy^*=y^*y=(1-\lambda)^21_\K$.
	
	The element $g=y\oplus y\in\B(\H)$ commutes with both $v_1$ and $v_2$,
	and so $y=\mu 1_\K$, for some $\mu\in\mathbb C$ such that $|\mu|^2=1-\lambda^2$. 
	Thus,
	\[
	v_2=\left[\begin{array}{cc} \lambda 1_\K & e^{i\theta}\sqrt{1-\lambda^2} 1_\K \\ e^{-i\theta}\sqrt{1-\lambda^2}1_\K & -\lambda 1_\K\end{array}\right],
	\]
	for some $\theta\in\mathbb R$.
As $v_1$ and $v_2$ are jointly unitarily equivalent to
	\[
	\bigoplus_1^{\mbox{\textrm dim}\,\K}\left[ \begin{array}{cc}-1 &0 \\ 0 & 1\end{array}\right]
	\,\mbox{ and }\,
	\bigoplus_1^{\mbox{\textrm dim}\,\K}
	\left[ \begin{array}{cc}\lambda &e^{i\theta}\sqrt{1-\lambda^2} \\ e^{-i\theta}\sqrt{1-\lambda^2} & -\lambda\end{array}\right],
	\]
	the only way for $\{v_1,v_2\}$ to be an irreducible set is for $\mbox{\textrm dim}\,\K=1$.
	Hence, $\H=\K\oplus\K$ has dimension $2$.
\end{proof}

\end{document}
